\chapter{Stellar Evolution}\label{ch:v}

We know that stars are born into molecular clouds of atomic Hydrogen, Helium and so on. But for stars to ignite, molecular clouds have to first cool down, for example by radiating away the excess heat.
\par It's clear that in molecular clouds there's a complex interplay between cooling and heating mechanisms. The latter, for example, are powered by the passing of cosmic rays that ionize the atoms of the cloud, fast electrons that bump against other molecules and atoms thus heating them up. Or maybe there are some nearby stars that are emitting X-rays, or high energy photons in general, that are able to heat up the cloud.
\par To cool down the cloud is forced to emit in the IR since it is optically thick to the higher frequencies (optical, UV and so on). Emission in the IR range occurs mainly by the de-excitation of collisionally excited molecules returning to the ground state. Similarly, the cloud can cool down through radiative emission of dust, which is heated up by nearby stars and can be assumed to emit as a black body at a temperature of roughly $1000\unit{K}$\footnote{Dust cannot be emitting at higher temperatures simply because it starts melting.}.
\par If the cloud weren't cooling down, it would be kinda a problem since no stars would be able to form. In fact, consider the Jeans' mass 
\begin{equation*}
    M_{J} \propto \left(\frac{T^{3}}{\rho}\right)^{1/2}
\end{equation*} 
which is the minimum mass for a gravitationally bound system to be unstable agaisnt gravitational collapse. If the molecular cloud were to be shrinking adiabatically, it is possible to show that the Jeans' mass would increase with the density, thus leading to a monolithic collapse that would kinda suck for the purpose of forming non-popIII stars.
\par What happens is exactly the opposite. At first the cloud is in hydrostatic equilibrium\footnote{We're neglecting the effects of rotations, magnetic fields etc. that would introduce a further complication that is beyond the aim of the present discussion.}, but then, due to some density perturbation, the cloud starts to shrink around its colder regions. This might happen for a whole series of reasons, for example the explosion of a nearby SN whose shockwave temporarily compresses part of the cloud inducing a collapse.
\par In this first phase the collpase is isothermal: The cloud is so sparse and the density so low ($n\approx 10-1000\unit{cm}^{-3}$) that the effect of opacity is essentially negligible, and thus the energy acquired through the gravitational collapse is immediately lost.
\par This thing that is shrinking is not a star, nor a protostar if we go by the strict meaning of the term\footnote{A protostar would be a star as it grows up, when it is still accreting matter from the cloud.}, it is but a slightly more dense than average blob of matter. In this blob, density quickly increases as the blob shrinks, but the temperature stays roughly the same, hence the Jeans' mass decreases. The cloud incurs into fragmentation.
\par As the blob shrinks, the inner regions become denser and denser, and only when it reaches density of order $10^{-13}\unit{g}\unit{cm}^{-3}$ the blob becomes slightly opaque. Therefore it starts to retain part of the gravitational energy that would otherwise be depleted contracting adiabatically. It is in that moment that a first semblance of a core forms.
\par This first core is a conglomerate of matter (of order of $0.01M_{\odot}$) in hydrostatic equilibrium and it's called \emph{first Larson's core}. On this clump of matter, there's a continuous infalling of other matter that leads to the formation of a shock front, where matter passes from supersonic to subsonic. In this region the kinetic energy of the infalling matter is released.
\par This first core, although in hydrostatic equilibrium, is unstable: As temperature increases, molecular Hydrogen dissociates ($\sim 2000\unit{K}$) and a new isothermal contraction begins: The excess energy released by the gravitational collapse is used to dissociate Hydrogen, thus leading to an expanding dissociation front from the hotter inner regions to the outskirts of the clump.
\par Now, this configuration is clearly unstable, for a little perturbation in density cannot still be matched by an analogous fluctuation in the pressure (remember that, at this stage, temperature isn't changing). Thus hydrostatic equilibrium is lost and the proto-core shrinks adiabatically once more, getting hotter and denser. 
\par This process continues until the protostar is born and a \emph{second Larson's core} is formed. This new core is in hydrostatic equilibrium and has a well-distinguishable structure for the rest of the cloud. The shock front is now on the surface of the core and matter is accreted with basically zero kinetic energy. 
\par Note that although it can be considered in hydrostatic equilibrium, it hasn't yet reached the conditions for thermal equilibrium. However, if we consider only what's within the shock front, we might as well be considereing that as a proper star.
\par At this stage the second Larson's core is about $1-50M_{\text{Jov}}$, so well below the mass of the Sun, while the radius is still rather large, ranging from $1-10R_{\odot}$, but there are still several criticalities on its determination, first and foremost the fact that the accretion rate is generally not known.
\par Consider that in the literature you might find values of the accretion rate $\dot{M}$ ranging from $10^{-7}-10^{-3}M_{\odot}\unit{yr}^{-1}$, where the latter is for burst-type accretions. On top of that, we don't even know the temporal behavior of $\dot{M}$, so we can't really tell if the accretion is continuous, decreasing with time or is a localized burst.
\par We don't know how the accretion rate $\dot{M}$ varies with time but, moreover, we can't even see the star, since it's all nested within the cloud. Around the accreting star therre is a region called \emph{opacity gap}, where atomic Hydrogen gets dissociated by the energy released from the shock front near the core.
\par Note that atomic Hydrogen is essentially transparent, so in this region opacity is particularly low. Out of the opacity gap there's the rest of the cloud, not yet dissociated.
\par Of course, we don't even know the geometry that is followed by the accretion process: It is believed that the accretion must be necessarily spherical at first, only to later settle into accretion from a disk: The star's wind and radiation pressure sweep away part of the cloud, leaving only a thick disk from which the star continues to accrete.
\par At the end of the star's accretion process, what's left of the disk will be very thin and many believe that's how a protoplanetary disk is formed.
\par Despite this pretty pessimistic introduction there are also some good news. When the star reaches its total mass, its characteristics are essentially determined, apparently forgetting whatever the properties of the protostar were before that moment. This is great news: Apart from the first $10^{5}-10^{6}\unit{yr}$ the star uses to accrete matter, stars all arrange themselves roughly in the same place in a HR diagram with similar characteristics.
\par This is what usually happens unless the accretion conditions are particularly extreme. Another quantity that determines the evolution during the accretion stage is the fraction of kinetic energy that is transferred to the star from infalling matter. We'll call the fraction that is transferred to the hydrostatic core $\alpha_{\text{acc}}$.
\par Now if $\alpha_{\text{acc}}\gtrsim 14\%$, the accretion is \emph{hot}, otherwise it's \emph{cold}. Stars with cold accretion have extreme accretion conditions: During the accretion, the protostar shrinks too much, leading to strikingly different characteristics from other stars that have formed normally.
\par This characteristics, however, are in opposition to experimental observations, so some believe cold accretion isn't really a thing.
\par We set the \emph{birth line}--so when a protostar becomes effectively a star--at the moment when the protostar's accretion rate gets below $10^{-7}M_{\odot}\unit{yr}^{-1}$. After the birth line, the star ufficially enters the \emph{pre-main sequence} stage.
\par As we've perhaps mentioned in the last chapter, these stars are still quite cold, so there isn't much they can burn. The first element they will ignite is Deuterium, at a temperature of $\sim 1\power{10}{6}\unit{K}$. This condition can be achieved both in the protostar phase and the pre-main sequence stage.
\par If Deuterium is ignited in the protostar stage, clearly it won't be stopping any time soon, since matter is still continuously falling into the star. This process continues until the star reaches its total mass and the molecular cloud is completely swept away.
\par If, instead, Deuterium is ignited during pre-main sequence stage, once it is completely depleted, the burning stops.
\par When a protostar "turns" into a star, it is still quite bloated, since during the accretion phase they have heated up and consequently expanded. But, they're still rather cold and thus quite luminous because they're all bloated up.
\par Since they're cold, the opacity is rather high. We know that opacity increases as the temperature decreases because of all the energy jumps, photoionization levels and so on. If the opacity is hight, the radiative gradient is larger than the adiabatic gradient and the star is, most of the times, completely convective. It might have a small radiative core in the inner, hotter regions, but as soon as Deuterium is ignited it vanishes, because something else happen at that point.
\par We know that the radiative gradient is proportional to the energy flux in cases like these. The star is convective because opacity is high. When Deuterium is ignited, the star heats up and the flux increases. The radiative gradient is still larger than the adiabatic one, so even if it has entered the pre-main sequence stage with a small radiative core, as soon as Deuterium is ignited it is lost.
\par The burning of Deuterium hinders the contraction of the star for a certain time and allows the star to "forget" everything that had happened before that moment. In the initial phase the star is completely mixed up so, even if at first it was dishomogeneous after all the complex accretion mechanisms, this dishomogeneity is lost.
\par One way or another, stars always pass through a stage where they are completely convective.

\section{Pre-Main Sequence}

We've mentioned somewhere in the "few" introducing lines of this chapter that during the formation of a star the yet-to-be protostar accretes matter from the surrounding molecular cloud. Such accretion might be distinguished into hot accretion or cold accretion (which seems to be against experimental observations). Up until accretion stops, the luminosity of the star can be thought as made up of two terms
\begin{equation*}
    L_{\text{tot}} = L_{\text{acc}}+L_{\text{sup}}
\end{equation*}
where $L_{\text{acc}}$ is the energy released through the accretion mechanism. The end of the protostar stage can be identified with the moment when $L_{\text{tot}} \approx L_{\text{sup}}$ and the accretion luminosity turns negligible. The speed of infalling matter is essentially equal to the escape velocity from the surface
\begin{equation*}
    0 = \frac{1}{2}mv_{f}^{2}-\frac{GM_{*}m}{R_{*}} \implies v_{f}\sim\sqrt{\frac{2GM_{*}}{R_{*}}}
\end{equation*}
hence the kinetic energy associated with a clump of infalling matter $\Delta m$ is just $E_{\text{k}} = \Delta m v_{f}^{2}/2$. The luminosity due to the accretion mechanism is just the variation over time of that kinetic energy (corrected with the efficiency of accretion $\alpha_{\text{acc}}$)
\begin{equation*}
    L_{\text{acc}} = (1-\alpha_{\text{acc}})\der{E_{k}}{t} = (1-\alpha_{\text{acc}})\frac{G\dot{M}M_{*}}{R_{*}}
\end{equation*}
and a typical relation is given
\begin{equation*}
    L_{\text{acc}} = 61L_{\odot}\left(\frac{\dot{M}}{10^{-5}M_{\odot}\unit{yr}^{-1}}\right)\left(\frac{M_{*}}{M_{\odot}}\right)\left(\frac{R_{*}}{5R_{\odot}}\right)^{-1}
\end{equation*}
As we've mentioned, at this stage the star is very bloated and cold, but its luminosity is rather high since 
\begin{equation*}
    L = 4\pi R^{2}\sigma T_{\text{eff}}^{4}
\end{equation*}
where the large extension compensates for the low temperature. Since the star is cold, it is also very opaque, hence the radiative gradient is 
\begin{equation*}
    \nabla_{\text{rad}} > \nabla_{\text{ad}}
\end{equation*}
and the star is basically everywhere convective. This is great news, because it means that the chemical composition of the star will be fairly homogeneous. On top of that, we know that, accretion or not, the newly born star will always be placed in the same spot in a HR diagram.
\par This is because the accreting stage is carried out in hydrostatic and (almost) thermal equilibrium and the star will be regardless reaching a total mass and complete convective-ness regardless of its previous accreting history (assuming it is not burst-like).
\par As our first model to describe a pre-main sequence star, we might avoid using the energy production equation (\ref{eq:energyproduction}) and assume adiabatic gradient everywhere but in the external regions, where superadiabaticity is not negligible\footnote{If we want to find an analytical model to describe a pre-main sequence star, we have to assume a constant adiabatic gradient. For example, assuming a monoatomic perfect gas that is also completely ionized, we'd find $\nabla_{\text{ad}} \sim 0.4$.}. At this stage, when nuclear reactions are still far from igniting, the evolution of the star is led by variations of the central temperature, that is highly influenced by the energy released during the gravitational collapse.
\par In his first works about the early stages of star evolution, Hayashi makes the following assumption
\begin{equation*}
    \der{\ln T}{\ln p} = \text{const.} = A
\end{equation*}
and then varies $A$. What he observed is that increasing $A$, the effective temperature decreases. This is because, if we assume negligible superadiabaticity, the largest gradient is the adiabatic one, so
\begin{equation*}
    \nabla_{\text{amb}} \approx \nabla_{\text{ad}}
\end{equation*}
Thus, Hayashi proved that fixed the mass and chemical composition of a star there exists a forbidden region for stars on an HR diagram, that is the region to the right of the line described by 
\begin{equation*}
     \der{\ln T}{\ln p} \sim 0.4
\end{equation*}
In the forbidden region, it can be shown that a star must have a temperature gradient greater than that limit. If a star is placed in the forbidden zone, with a temperature gradient much greater than 0.4, it will experience rapid convection that brings the gradient down. Since this convection will drastically change the star's pressure and temperature distribution, the star is not in hydrostatic equilibrium, and will contract until it is. Eventually, the star is brought (literally) back on tracks.
\par We thus define \emph{Hayashi's track} that limiting line along which all pre-MS stars will align, since they're both convective and in hydrostatic equilibrium. One striking feature of Hayashi's track is that, on an HR diagaram, it is almost vertical, meaning the star decreases its luminosity but keeps its effective temperature roughly constant all throughout its evolution.
\par This is because, given the low effective temperatures, the main source of opacity is that caused by the H$^{-}$ (\ref{eq:hydrogenopacity}). When the star shrinks, the temperature raises and so does the electron density and, consequently, the opacity. The radiative gradient quickly climbs up as well and so the change in temperature never really gets to the surface.
\par The evolution along the track is led by gravitational forces and so the timescale for the evolution is the Kelvin-Helmoltz
\begin{equation*}
    \tau_{KH} \propto \frac{GM^{2}}{RL}
\end{equation*}
but since the mass is not changing anymore it is effectively dependent only on the luminosity and the radius. Furthermore
\begin{equation*}
    L = 4\pi R^{2}\sigma T_{e}^{4}
\end{equation*}
and since the effective temperature is constant along the track $R \sim L^{1/2}$ meaning that $\tau_{KH}\propto L^{-3/2}$. Hence the luminosity of the star is decreasing with time
\begin{equation*}
    L \propto t^{-2/3}
\end{equation*}
which is consistent with the star shrinking its radius along the track.
\subsection{Analytical model}
Although not particularly enticing, it is possible to find an analytical description for pre-MS stars. First of all we might describe the inner regions with a politropic equation
\begin{equation*}
    p = K_{2}\rho^{\gamma} \quad \gamma = \frac{1}{1-\nabla_{\text{ad}}}
\end{equation*}
which we might easily integrate as we've seen in earlier chapters. This, however, wouldn't work that well if we were to extend it all up to the atmosphere, because that's where our adiabatic convection approximation breaks down. A reasonable approximation to deal with the more external layers of a stellar atmosphere would be assuming radiative transport, using, say, the grey atomsphere approximation (\ref{eq:greyatmosphere_temp}).
\par Even if the convection within the star reaches layer above $\tau = 2/3$ (which is the demarcation line between the inner regions and the atmosphere in the grey atomsphere approximation), we can always perform the atmospheric integration down to the point where convection sets in and take the boundary conditions there.
\par If we assume 
\begin{equation*}
    \der{p}{\tau} = \frac{GM}{\bar{\kappa}R^{2}}
\end{equation*}
where the constant value of the opacity through the atmosphere is equal to $\bar{\kappa} = \kappa_{0}p^{a}T^{b}$ (which is analogue to Cramer's). It is possible to show, see for example \cite{2005essp.book.....S} §4.3.1, that if we link on a given boundary the solutions obtained by integrating the politropic equation and integrating over the atmosphere layers, we'd obtain a relation of this type 
\begin{equation}
    \log(T_{\text{eff}}) = A\log\left(\frac{L}{L_{\odot}}\right)+B\log\left(\frac{M}{M_{\odot}}\right)+\text{const}
    \label{eq:luminosity_pre-ms}
\end{equation}
where $A$ and $B$ are two constants
\begin{equation*}
    A = \frac{0.75a-0.25}{b+5.5a+1.5} \quad B = \frac{1.5a-0.5}{b+5.5a+1.5}
\end{equation*}
Using realistic values for which the main source of opacity is the H$^{-}$, you'd obtain $A\sim 0.01$ and $B\sim 0.14$. The slope of the Hayashi track in the HR-d is therefore 
\begin{equation*}
    \parder{\log L}{\log T_{\text{eff}}} \sim \frac{1}{A} > 10
\end{equation*}
meaning the Hayashi track is basically vertical; also, since
\begin{equation*}
    \parder{\log T_{\text{eff}}}{\log M} \sim 0.14
\end{equation*}
its location has a mild dependence on the stellar mass, and it moves towards higher $T_{\text{eff}}$ for increasing masses.
\par As soon as the star climbs down the Hayashi track, Deuterium is immediately ignited\footnote{The required temperature, as we've seen, is of $T\sim 10^{6}\unit{K}$, has a $Q$ value of $\sim 5.5\unit{MeV}$ and an initial abundance $X_{D}\sim 2\power{10}{-5}$.}
\begin{equation*}
    \element{D}{2}{1}+\element{p}{1}{1} \to \element{He}{3}{2}+\gamma
\end{equation*}
The timescale for burning Deuterium is given by 
\begin{equation*}
    \tau_{D} = \frac{1}{n_{D}\langle \sigma v \rangle_{\text{D}}+p}\sim 10^{6}\unit{yr}
\end{equation*}
and the energy production rate is $\epsilon_{D}\propto T^{11-12}$. Note that if the burning of Deuterium is ignited in the protostar stage, this will go on as long as accretion feeds fresh Deuterium to the star.
\par It is also possible to show that an increase in metallicity (which causes an increase of the opacity) shifts the track to lower effective temperatures, whereas an increase of Helium at constant metallicity has the opposite effect. This can be roughly explained recalling the EoS for a perfect gas 
\begin{equation*}
    p = \frac{\rho}{\mu m_{\text{H}}}kT
\end{equation*}
If we assume the pressure to be fixed by the request of the star to be in hydrostatic equilibrium, an increase in Helium abundance increases the mean molecular weight and hence a higher temperature is needed.
\par Meanwhile, in the external regions, the effect of convection is more relevant. In fact, if the mixing length coefficient $\alpha_{ML}$ decreases, the convective flux decreases (and the radiative one increases). Necessarily, the ambient temperature gradient has to ramp up and hence the effective temperature will decrease.
\par When the star finally becomes more radiative, it strays aways from the Hayashi track, and enters the so-called Henyey stage. Here nuclear reactions have yet to begin, and stars present different trends for luminosity depending on their mass
\begin{equation*}
    L_{\text{LMS}}\propto T_{\text{eff}}^{0.8}M^{4.4} \quad L_{\text{UMS}}\propto M^{3}
\end{equation*}
Less massive stars (VLMS, $M\lesssim 0.3M_{\odot}$) have other problems seems they keep being convective and therefore stick to the Hayashi track\footnote{Note that for stars in the range $M \lesssim 0.4M_{\odot}$ the $\element{He}{3}{2}$ becomes a pseudo-primary element, never reaching equilibrium.}.
\par We define \emph{zero age main sequence} the beginning of the MS, which has to be identified with the moment during which the star is entirely supported by Hydrogen burning and all secondary elements have reached equilibrium\footnote{Note that this might need a lot of time depending on the reaction cross sections. In the CN-NO bi-cycle, the main postponer of the beginning of MS is $\element{N}{14}{7}$, which has a tiny cross section.}.
